\section{Binary generators and coset representatives}\label{app:codes}

Table~\ref{tab:generators} gives the generators used for the three
dense sign codes. An entry is a hexadecimal mask, representing the
binary word $c$ by $\sum_i c_i2^i$. The generator of length $32$ is
used unchanged in ambient dimensions $32$--$37$. The
length-$38$ and length-$39$ columns are the kernels reconstructed
from~\eqref{eq:qr}. Table~\ref{tab:representatives} supplies the affine
representatives; the length-$32$ code uses only zero.

\begin{table}[ht]
\centering\small
\caption{Generator masks for $H_{32}$, $H_{38}$, and $H_{39}$. Blank
entries reflect their different dimensions.}
\label{tab:generators}
\begin{tabular}{rrrr}
\toprule
Row & $H_{32}$ & $H_{38}$ & $H_{39}$\\
\midrule
0 & \texttt{4EF0C081} & \texttt{2000C0DE0B} & \texttt{4000152727}\\
1 & \texttt{24A09082} & \texttt{1000AA229D} & \texttt{2000C0DE0B}\\
2 & \texttt{48609084} & \texttt{8009F5CD6} & \texttt{1000AA229D}\\
3 & \texttt{1E505088} & \texttt{4004FAE6B} & \texttt{8009F5CD6}\\
4 & \texttt{6C905090} & \texttt{200ED9AAD} & \texttt{4004FAE6B}\\
5 & \texttt{FA30C0A0} & \texttt{100BC80CE} & \texttt{200ED9AAD}\\
6 & \texttt{AAC000C0} & \texttt{805E4067} & \texttt{100BC80CE}\\
7 & \texttt{D2608100} & \texttt{40E56DAB} & \texttt{805E4067}\\
8 & \texttt{7E60D200} & \texttt{20B8FB4D} & \texttt{40E56DAB}\\
9 & \texttt{B8601400} & \texttt{1096303E} & \texttt{20B8FB4D}\\
10 & \texttt{D800D800} & \texttt{84B181F} & \texttt{1096303E}\\
11 & \texttt{CC606000} & \texttt{4EFC197} & \texttt{84B181F}\\
12 & \texttt{1EE10000} & \texttt{2BDAD53} & \texttt{4EFC197}\\
13 & \texttt{B2B20000} & \texttt{1949B31} & \texttt{2BDAD53}\\
14 & \texttt{74740000} &  & \texttt{1949B31}\\
15 & \texttt{D8D80000} &  & \\
16 & \texttt{FF000000} &  & \\
\bottomrule
\end{tabular}

\end{table}

\begin{table}[ht]
\centering\small
\caption{Representatives of the eleven cosets in length $38$ and the
ten cosets in length $39$.}
\label{tab:representatives}
\begin{tabular}{rrr}
\toprule
$a$ & $r_a$ in length 38 & $r_a$ in length 39\\
\midrule
0 & \texttt{0} & \texttt{0}\\
1 & \texttt{C6E94} & \texttt{ADCFB}\\
2 & \texttt{DF6ABF} & \texttt{F7795F}\\
3 & \texttt{27B441} & \texttt{740E29}\\
4 & \texttt{D3CE66} & \texttt{940DFF}\\
5 & \texttt{41BBBB} & \texttt{3B66AB}\\
6 & \texttt{BC1E1F} & \texttt{6CD301}\\
7 & \texttt{52727} & \texttt{4709D4}\\
8 & \texttt{9951C8} & \texttt{98A926}\\
9 & \texttt{7001EB} & \texttt{D23688}\\
10 & \texttt{3F6969} & \\
\bottomrule
\end{tabular}

\end{table}

The complete kernel weight distributions are given in
Table~\ref{tab:weights}. The least nonzero weights give the within-coset
distance bounds. For every distinct pair of representatives in
Table~\ref{tab:representatives}, the affine space
$r_a+r_b+H_m$ has minimum weight ten. Enumeration of these spaces
establishes both the cross-coset distances and disjointness of the
cosets. This determines each code in Proposition~\ref{prop:signs}
entirely from the printed data.

\begin{table}[ht]
\centering\small
\caption{Weight distributions $A_w=|\{h\in H_m:\wt(h)=w\}|$.
Every omitted weight has coefficient zero.}
\label{tab:weights}
\begin{tabular}{rrrr}
\toprule
$w$ & $H_{32}$ & $H_{38}$ & $H_{39}$\\
\midrule
0 & 1 & 1 & 1\\
8 & 908 & 0 & 0\\
10 & 3328 & 0 & 0\\
12 & 14784 & 669 & 970\\
14 & 27392 & 0 & 0\\
16 & 38246 & 5289 & 8952\\
18 & 27392 & 0 & 0\\
20 & 14784 & 8002 & 16452\\
22 & 3328 & 0 & 0\\
24 & 908 & 2310 & 5988\\
28 & 0 & 113 & 402\\
32 & 1 & 0 & 3\\
\bottomrule
\end{tabular}

\end{table}
